
\documentclass{article}
\usepackage{colortbl}
\usepackage{makecell}
\usepackage{multirow}
\usepackage{supertabular}

\begin{document}

\newcounter{utterance}

\centering \large Interaction Transcript for game `cladder', experiment `full\_v1.5\_default', episode 4315 with Qwen3.5{-}4B.
\vspace{24pt}

{ \footnotesize  \setcounter{utterance}{1}
\setlength{\tabcolsep}{0pt}
\begin{supertabular}{c@{$\;$}|p{.15\linewidth}@{}p{.15\linewidth}p{.15\linewidth}p{.15\linewidth}p{.15\linewidth}p{.15\linewidth}}
    \# & \multicolumn{2}{c}{Player} && \multicolumn{2}{c}{Game Master} \\
    \hline

    \theutterance \stepcounter{utterance}  
    & & & \multicolumn{4}{p{0.6\linewidth}}{
        \cellcolor[rgb]{0.9,0.9,0.9}{
            \makecell[{{p{\linewidth}}}]{
                \texttt{\tiny{[P1$\langle$GM]}}
                \texttt{You are an expert in causal inference. The following question is not a typical commonsense query, but rather a meticulously designed question created by a professor specializing in causal inference, intended to assess the students' mastery of the course content.} \\
\texttt{Imagine a self{-}contained, hypothetical world with only the following conditions, and without any unmentioned factors or causal relationships: Maternal smoking status has a direct effect on infant's birth weight and freckles. Health condition has a direct effect on infant's birth weight and freckles. Infant's birth weight has a direct effect on freckles. Health condition is unobserved. For infants with nonsmoking mothers and low infant birth weight, the probability of freckles is 46\%. For infants with nonsmoking mothers and normal infant birth weight, the probability of freckles is 21\%. For infants with smoking mothers and low infant birth weight, the probability of freckles is 74\%. For infants with smoking mothers and normal infant birth weight, the probability of freckles is 47\%. For infants with nonsmoking mothers and with poor health, the probability of normal infant birth weight is 51\%. For infants with nonsmoking mothers and with good health, the probability of normal infant birth weight is 72\%. For infants with smoking mothers and with poor health, the probability of normal infant birth weight is 3\%. For infants with smoking mothers and with good health, the probability of normal infant birth weight is 35\%. The overall probability of good health is 90\%. If we disregard the mediation effect through infant's birth weight, would maternal smoking status negatively affect freckles?} \\
\texttt{Start your answer with 'Yes' or 'No', followed by additional reasoning or evidence to support your explanation.} \\
            }
        }
    }
     \\ \\

    \theutterance \stepcounter{utterance}  
    & \multicolumn{4}{p{0.6\linewidth}}{
        \cellcolor[rgb]{0.561,0.737,0.561}{
            \makecell[{{p{\linewidth}}}]{
                \texttt{\tiny{[P1$\rangle$GM]}}
                \texttt{Yes To determine whether maternal smoking status negatively affects freckles while disregarding the mediation effect through infant's birth weight, we must calculate the Total Effect of smoking on freckles. This total effect is the sum of the direct effect of smoking on freckles and the indirect effect (which is zero in this specific scenario because we are instructed to disregard the mediation path). We can derive the Total Effect by comparing the marginal probability of freckles between smokers and non{-}smokers, accounting for the distribution of health conditions. Let S be smoking status (1 = Smoking, 0 = Non{-}smoking), H be health status (1 = Good, 0 = Poor), W be birth weight (1 = Normal, 0 = Low), and F be freckles. Given P(H=1) = 0.90, therefore P(H=0) = 0.10. For Non{-}smokers (S=0): P(W=1 | S=0, H=1) = 0.72, P(W=1 | S=0, H=0) = 0.51, P(F=1 | S=0, W=1) = 0.21, P(F=1 | S=0, W=0) = 0.46. For Smokers (S=1): P(W=1 | S=1, H=1) = 0.35, P(W=1 | S=1, H=0) = 0.03, P(F=1 | S=1, W=1) = 0.47, P(F=1 | S=1, W=0) = 0.74. We calculate the Marginal Probability of Freckles for Non{-}Smokers P(F=1 | S=0) using the law of total probability, conditioning on Health (H) and Birth Weight (W). The calculation is: P(F=1 | S=0) = P(F=1 | S=0, W=1, H=1) * P(W=1 | S=0, H=1) * P(H=1) + P(F=1 | S=0, W=0, H=1) * P(W=0 | S=0, H=1) * P(H=1) + P(F=1 | S=0, W=1, H=0) * P(W=1 | S=0, H=0) * P(H=0) + P(F=1 | S=0, W=0, H=0) * P(W=0 | S=0, H=0) * P(H=0). Substituting the values: P(F=1 | S=0) = 0.21 * 0.72 * 0.9} \\
            }
        }
    }
    & & \\ \\

    \theutterance \stepcounter{utterance}  
    & & & \multicolumn{2}{p{0.3\linewidth}}{
        \cellcolor[rgb]{0.9,0.9,0.9}{
            \makecell[{{p{\linewidth}}}]{
                \texttt{\tiny{[GM$|$GM]}}
                \texttt{Yes} \\
            }
        }
    }
    & & \\ \\

    \theutterance \stepcounter{utterance}  
    & & & \multicolumn{2}{p{0.3\linewidth}}{
        \cellcolor[rgb]{0.9,0.9,0.9}{
            \makecell[{{p{\linewidth}}}]{
                \texttt{\tiny{[GM$|$GM]}}
                \texttt{no} \\
            }
        }
    }
    & & \\ \\

    \theutterance \stepcounter{utterance}  
    & & & \multicolumn{2}{p{0.3\linewidth}}{
        \cellcolor[rgb]{0.9,0.9,0.9}{
            \makecell[{{p{\linewidth}}}]{
                \texttt{\tiny{[GM$|$GM]}}
                \texttt{game\_result = LOSE} \\
            }
        }
    }
    & & \\ \\

\end{supertabular}
}

\end{document}
